\newpage
\section{DISCUSSION}
%%%%%%%%%%%%%%%%%%%%%%%%%%%%%%%%%%%%%%%
%% Discuss your thesis here

%\renewcommand{\thefigure}{\thesection.\arabic{figure}} %setzt Nummerierung der Figure wieder auf 0 für Section
%\setcounter{figure}{0}

%\renewcommand{\thetable}{\thesection.\arabic{table}} %setzt Nummerierung der Table wieder auf 0 für Section
%\setcounter{table}{0}

\subsection{Biochemistry}



\subsubsection{General characterization of \textup{Methanobrevibacter}}


General microbiological characterization experiments were performed with \acrshort{thau}, \acrshort{bovi}, \acrshort{milli}, and \acrshort{woli} to improve our understanding of the poorly understood \acrlong{mbrevi}, which is the dominant genus of methanogenic \gls{archa} in the rumen of cattle \parencite{baca-gonzalezAreVaccinesSolution2020}. \\

\textbf{Cultivation} 

\noindent In this bachelor thesis, the \acrlong{mbrevi} species were successfully cultivated in liquid and solid DSM~1523 and 119. Unlike typical cultivation media described in the literature, these media do not contain rumenfluid \parencite{liMethanobrevibacterBoviskoreaniJH1T2023, leeMethanobrevibacterBoviskoreaniSp2013, millerDescriptionMethanobrevibacterGottschalkii2002}. Most of the experiments were performed in DSM 1523, because there was no visible difference in cultivation between the two media, and DSM~1523 is easier to handle since it does not contain the toxic fatty acid mixture and the sludge fluid. In general, \acrlong{mbrevi} reach their growth maximum within two to four days, with a maximum \acrshort{OD} of 0.3 that does not exceed 0.4. Further growth experiments with a higher time resolution are needed to calculate general growth parameters such as generation time or growth rate, particularly during the exponential growth phase. The highest resolution achieved here was once per day. The detected growth is as fast as that described in the literature for other \gls{archa}. However, in this bachelor thesis, the \acrlong{mbrevi} did not grow as densely, in some cases reaching only about one-third of the optical density reported by \textcite{liMethanobrevibacterBoviskoreaniJH1T2023, hoedtDifferencesDownunderAlcoholfueled2016}.

A closer look at \acrshort{bovi} reveals that it grows reliably in the described manner and can also be easily spread plated. The only unusual aspect is the formation of precipitates (\cref{fig:bovis_präzipitat}), which makes photometric measurements somewhat more difficult. In this case, it is important to homogenize the culture as thoroughly as possible beforehand, even if the precipitates cannot be completely dissolved. Interestingly though, this behavior is not described for the type strain (JH1\textsuperscript{T}) of this species and also the other investigated \acrshort{mbrevi} species did not show this. Why \acrlong{bovi} precipitates remains unclear.  

\acrshort{thau}, on the other hand, grew much more unreliably, as greater heterogeneity in the timing of growth onset and growth density was observed both within triplicates and between experiments. In addition, \acrshort{thau} lysed quickly after reaching maximal growth, such that within 4 to 6 days the culture density was low enough to appear visually clear of cells. Furthermore, spread plating on DSM~1523 (1.5~\% agar) and pour \gls{plateing} with DSM~119 (1~\% agar) did not work (\cref{fig:spread_plating} and \cref{fig:pour_plating}). Roll plating was successful, but with a notably lower \gls{plateing} efficiency than for \acrshort{bovi} and \acrshort{woli} (\cref{tab:roll_plat_eff}).
This growth behavior of \acrshort{thau} makes it more difficult to analyze. Thus, \acrshort{thau} does not appear to be suitable for basic research due to its unreliable growth.

Future research should further investigate whether \acrshort{bovi} and \acrshort{thau} can use a liquid substrate as an alternative to \ch{H2}, potentially allowing the upscaling of growth experiments through the use of a plate reader. In contrast to the literature -- where \textcite{liMethanobrevibacterBoviskoreaniJH1T2023} reported that \textit{M. boviskoreani} JH1\textsuperscript{T} could use \acrshort{etoh} as a substrate -- neither \acrshort{thau} nor \acrshort{bovi} was able to use \acrshort{etoh} as an alternative substrate here.
Another substrate that some \acrlong{mbrevi} species can use, according to the literature, is formate, which should therefore be tested next \parencite{liMethanobrevibacterBoviskoreaniJH1T2023, hoedtDifferencesDownunderAlcoholfueled2016}.\\

\textbf{Cellnumber per mL}

\noindent For the first time cellnumber per mL was determined for \acrshort{thau}, \acrshort{milli}, \acrshort{woli}, and \acrshort{bovi}. Compared to the cellnumber per mL for \acrshort{ecol} ($8 \cdot 10^8$~cells/mL) at a \acrshort{OD} of 1 given by \textcite{merckkgaaCellDensityMeasurement2026} the counted cells/mL for the \acrshort{mbrevi} cultures are quite high (e.g. \acrshort{woli}: $10.819 \cdot 10^8$~cells/mL at a \acrshort{OD} of 0.387; \cref{tab:ø_cell_per_culture}).  
Since only one or two cultures have been counted for these organisms so far, further investigation are needed to verify these and investigate the relationship between cell count and \acrshort{OD}.

Besides the small data sample, two further factors contribute to the increased uncertainty.
First, the chamber depth of 0.1~mm \parencite{lo-laboroptikgmbhZaehlkammer2026} caused cells to appear in multiple focal planes, increasing the risk of overlooked or double-counted cells.
Second, the still small size of the cells at a 400x magnification plus the tendency to form chains or clumps, increase the risk of overlooked or double-counted cells further due to the problem to distinguish the cells from one another (\cref{fig:cell_count_bovis_bsp}).

One approach to mitigate this risk was to count 150 to 250 cells per large square, rather than less than 150. Thereby, reducing the relative impact of individual miscounts. This, of course, led to an increased workload.
Further approaches could utilize the intrinsic fluorescence of methanogenic \gls{archa}  \parencite{lambrechtFlowCytometricQuantification2017} and perform the counting using fluorescence microscopy. This would make it easier to identify the small \acrlong{mbrevi} cells at 400x magnification.\\

\textbf{Roll plating}

\noindent In general, the technique works well for \acrshort{bovi}, \acrshort{thau} and \acrshort{woli}. For \acrshort{bovi} and \acrshort{woli} less starting material should be used next time, since there were too many colonies, which are consequently too small for selective picking (see \cref{fig:roll_platt} and \cref{tab:roll_plat_eff}). Therefore, it is not proven through sequencing that the right \acrlong{mbrevi} grew, but the constant consumption of the \ch{H2}/\ch{CO2} atmosphere proofs there was definitely anaerobic growth.
Still this plating technic has the potential to be a time saving alternative to the spread plating due to less work inside the \gls{glove}. However, it needs improvement for a fitting quantity, so that less but bigger colonies grow. 

The plating efficiencies shown in \cref{tab:roll_plat_eff} were very low, with values ranging from below 1 to 20~ppm. By contrast, \textcite{apolinarioPlateColonizationMethanococcus1996} reached an efficiency of over 50~\% for other \gls{archa} species. However, the calculated plating efficiencies in this bachelor thesis should be interpreted with caution, since the \acrshort{OD} of the initial solution was not measured. Instead, it was assumed retrospectively that the cultures had a similar OD to those obtained by cell counting.
For future plating efficiency determination, it is important to measure the \acrshort{OD} in advance and to plate only at an \acrshort{OD} corresponding to those in \cref{tab:cellnumb}.

%%%%%%%%%%%%%%%%%%%%%%%%%%%%%%%%%%%%%%%%%%%%%%%%%%%%%%%%%%%%%%%%%%%%%%%%%%%%%%

%%%%%%%%%%%%%%%%%%%%%%%%%%%%%%%%%%%%%%%%%%%%%%%%%%%%%%%%%%%%%%%%%%%%%%%%%%%%%%

%%%%%%%%%%%%%%%%%%%%%%%%%%%%%%%%%%%%%%%%%%%%%%%%%%%%%%%%%%%%%%%%%%%%%%%%%%%%%%

%%%%%%%%%%%%%%%%%%%%%%%%%%%%%%%%%%%%%%%%%%%%%%%%%%%%%%%%%%%%%%%%%%%%%%%%%%%%%%


\subsubsection{Sensitivity tests}

\textbf{Puromycin}

\noindent Since the \gls{selection} of successfully transformed \acrlong{mbrevi} is based on resistance to \acrfull{puro}, it was first necessary to determine whether, and at what concentration, \acrshort{bovi}, \acrshort{milli} and \acrshort{thau} are sensitive to this antibiotic.
\acrshort{thau} is sensitive to \acrshort{puro} at a concentration at 10~µg/mL and showed no spontaneous resistance within 10 days. \acrshort{bovi} is also sensitive at a concentration of 10~µg/mL, but 20~µg/mL enhances the time until spontaneous resistance occurs. However, at 20~µg/mL spontaneous resistance occurred after 8 days. Hence, cells have to be collected within this time span.

\Acrshort{puro} can be used for selection of transformed \acrshort{bovi} and \acrshort{thau}, although they are somewhat less sensitive than other \gls{archa} species. For example, \textit{Methanococcus voltae} has a \acrfull{mic} of 2~µg/mL, and \textit{Methanosarcina mazei} has an \acrshort{mic} of 5~µg/mL \parencite{possotAnalysisDrugResistance1988, mondorfNovelInducibleProtein2012}. For the upcoming conjugation experiments, a concentration of 20~µg/mL~\acrshort{puro} was chosen for selection. Should growth occur within 8~days, it can be assumed that this was mediated by the introduced resistance and not due to spontaneous resistance.\\

\textbf{\acrshort{azaeight}}

\noindent The \gls{markerless} mutagenesis is based on the \gls{counterselection} with \acrshort{azaeight}, therefore \acrshort{bovi}, \acrshort{thau} and \acrshort{milli} has to be sensitive towards the base analogue. Although the \textit{hpt} gene on the plasmid already provides a backup mechanism that induces artificial sensitivity to \acrshort{azaeight}, as suggested in \textcite{kleinMarkerlessMutagenesisEnables2025}. The sensitive experiment testing concentrations ranging from 100 to 1000~µg/mL~\acrshort{azaeight} have little to no informativeness, due to the discovery of a sensitivity towards the solvent of \acrshort{azaeight} \acrfull{dmso} (\cref{fig:aza_mbov} through \cref{fig:aza_mmill}).

However, for \acrshort{bovi} a sensitivity is detectable for concentration of 500~µg/mL based on a two days delayed growth start and around 0.7 lower maximal \acrshort{OD} values (\cref{fig:aza_mbov}). \acrshort{milli} growth results do not allow a clear conclusion due to heterogenous growing through the different concentration preparation as well as between the triplicates (\cref{fig:aza_mmill}). Nevertheless, the fact that all samples exhibit similar inhibition at the maximum concentration compared to the positive control may indicate that, even though the 100 µg/mL solution contains ten times less DMSO in the medium than the DMSO preparation, \acrshort{azaeight} also inhibits growth here. What argues against this, however, is that no stronger inhibition is observed at higher \acrshort{azaeight} concentrations, even though the \acrshort{dmso} concentration and base analog concentration increase. \acrshort{thau} is very sensitive against \acrshort{dmso} and showed no growth at all except in the positive control (\cref{fig:aza_mthau}).

A plausible explanation for the negative effect of DMSO on the growth of \acrshort{mbrevi} species is that, as reported by \citet{atcherEffectDMSOAqueous2013}, DMSO acts as a moderate oxidizing agent toward dithiols at room temperature. Therefore, \acrshort{dtt}, which is added to the culture as a reducing agent to remove residual oxygen and keep the medium anaerobic, reacts with the DMSO to form its cyclic structure \citep{sanzSimpleSelectiveOxidation2002}. Thereby, a disulfide bridge is formed, and \acrshort{dmso} is reduced to dimethyl sulfide. Consequently, no reducing agent remains present in the medium, and the anaerobic cultures are exposed to greater oxidative stress.

This hypothesis is supported by the following:
First, when samples were collected for photometric measurements, the cultures containing DMSO turned red more quickly and also took longer to lose their color than the positive controls(\cref{sec:azaeight_result}).
Second, in the \acrshort{thau} measurement series, one preparation contained only 1~mL \acrshort{dmso} and was not inoculated. Within a few days, this preparation culture turned distinctly red and did not fade (\cref{fig:dmso_neg_kontrolle}).
However, this also means that the \acrlong{mbrevi} themselves can partially reduce the solution, since in this case the cultures did not turn red as quickly or as distinctly as in the pure \acrshort{dmso} control, although they contained \acrshort{dmso} as well. Bute the \acrlong{mbrevi} do not appear to be able to fully compensate for the absence of \acrshort{dtt}, resulting in inhibited growth.

In the future, an alternative solvent should be sought or the counterselection can not be performed properly. In the literature \textcite{olsonTransformationClostridiumThermocellum2012} state 50~mg/mL \acrshort{azaeight} is dissolvable in 1~M \ch{NaOh} or \ch{KOH}. A possible problem here could be the influence on the pH of the media, the effect of this on the growth of the \acrlong{mbrevi} has to be investigated.  


%%%%%%%%%%%%%%%%%%%%%%%%%%%%%%%%%%%%%%%%%%%%%%%%%%%%%%%%%%%%%%%%%%%%%%%%%%%%%%

%%%%%%%%%%%%%%%%%%%%%%%%%%%%%%%%%%%%%%%%%%%%%%%%%%%%%%%%%%%%%%%%%%%%%%%%%%%%%%

%%%%%%%%%%%%%%%%%%%%%%%%%%%%%%%%%%%%%%%%%%%%%%%%%%%%%%%%%%%%%%%%%%%%%%%%%%%%%%


\subsubsection{Plasmid constructs}
\label{sec:plasmid_construction_diskussion}

An interesting discovery was made while constructing the plasmids for the transformation: the \acrlong{aa} sequence of the \textit{pac} gene, which was on the parent plasmid and originally planned to be on the new constructed ones for the \gls{conjug}, differs in 5 amino acids from the originally \acrlong{aa} sequence \textit{Streptomyces alboniger}. This resulted in a different structure prediction by \gls{alphafold} (\cref{fig:pac_org_and_old}), which, according to \textcite{kryshtafovychCriticalAssessmentMethods2021}, is expected to be as accurate as experimentally determined structures. Interestingly though, this differing sequence was used in the literature as the amino acid-sequence of \textit{Streptomyces alboniger} \parencite{lacalleMolecularAnalysisPac1989}.
If the changed protein-structure really has an influence on the function or the performed conjugation is not clear. To answer this question it needs further investigation. 
Just to be on the safe side, for the used plasmids in the \gls{conjug} \textit{pac} gene was replaced with the codon optimized original \textit{pac} from \textit{Streptomyces alboniger}.

In general the constructed plasmids followed the scheme in \cref{fig:general_plasmid_map}, where the selection of a successful conjugation works with a resistance against puromycin via \textit{pac}-gene and a counterselection with a sensitivity against \acrfull{azaeight}. 
If the whole-genome sequencing proofs the success of the transformation, a next step for the plasmid experiments would to test the other planed plasmids with a different promoter (\textit{PglnA} instead of \textit{PhmtB}, see \cref{tab:plasmid_konstrukte}).


\subsubsection{Transformation via conjugation}
\label{sec:discuss_trafo}

Promising initial results have been achieved for implementing the first genetic toolbox for \acrshort{bovi} and \acrshort{milli}, even though a final assessment by whole genome sequencing is still pending (see \cref{sec:transformation_via_conjugation}). The transformation of \acrshort{thau} was not successful.
Three cultures in two preparation lines for \acrshort{milli} and three cultures in one preparation line for \acrshort{bovi} (\cref{tab:conjugation_preparation}) were tested positive by PCR analysis to have taken up plasmid-DNA (\cref{fig:gelimage_conj_plasmid_check_epiderm}), although the fragment appears smaller than expected. 

This result has to be taken with caution, because the positive PCR signal could result from four other reasons.
First, the plasmid could still be present in the media and be extracted along with the \acrshort{gDna}. This scenario is maybe likely for the recovery culture but for the 1st or especially the 2nd transfer, where only 1~\% is left of the recovery culture, this seems to be unlikely. Second, the plasmid is amplified through \acrshort{ecol} through the corresponding \gls{ori} on the plasmid, but the bacterial \gls{S-sequenz} revealed no contamination with \acrshort{ecol} although \acrlong{sepidermidis} was detected (\cref{fig:gelimage_conj_plasmid_check_epiderm}). For this species the plasmid contains no \gls{ori}. Third, the primers have an off-target with a similar sized product in the gDNA of present microbes in the culture, which was disproved by a search using “Primer Blast” on NCBI. Fourth, it could be due to an unlikely contamination of the primer set or of a used solution with the plasmids, but this can not be ruled out due to a missing negative control. 

To have final clearance a whole genome sequencing is needed. However, whole genome sequencing  was not performed, because of the persisting contamination with. Because of that, the next step would be to purify these cultures through plating on selective media and picking a single colonie for the following analyses.
If whole-genome sequencing confirms that the plasmid has been successfully integrated, the next step would be to test the counter-selection implanted in the plasmid design using \acrshort{azaeight}.  Also, it would be interesting to test if other transformation methods such as electroporation could work.

Based on the comparison of the six successfully growing cultures, some factors influencing the \gls{conjug} can be assumed, although a rather small sample was investigated here and, to verify this, further investigation is definitely needed.

First, the growth status of \acrlong{mbrevi} seems to be relevant. Growth only occurred among the 16 of overall 28 preparation lines that used freshly grown cultures (\cref{tab:conjugation_preparation}). This aligns with the literature, where cultures from late exponential or early stationary phase were used for successful transformation \parencite{capovillaTransformationPlasmidDNA2025, finkShuttlevectorSystemAllows2021}. This could also be the reason why the transformation of \acrshort{thau} did not work, since only an old culture was used for this strain. The actual growth phase of the initial cultures used cannot be determined here, due to missing \acrshort{OD} values for the \acrlong{mbrevi} cultures. This should be addressed in future investigations.

Second, all three growing preparation lines used \acrshort{ecol26} as the donor strain for the interdomain \gls{conjug}. Future transformation attempts with \acrshort{bovi} and \acrshort{milli} should focus on this. The methylation pattern of \acrshort{ecol26} therefore does not impact the transformation of \acrshort{bovi} and \acrshort{milli} \parencite[SI]{schwalenBioinformaticExpansionDiscovery2018a}.

Third, the recovery time does not appear to be as critical for the success of the transformation as described for \textit{Methanothermobacter thermautotrophicus}~$\Delta$H in \citet{finkShuttlevectorSystemAllows2021}. However, a slight trend towards a shorter recovery time is visible, as 5 out of the 6 successfully grown cultures were among the preparation lines with a shorter recovery time. In addition, no direct trend is detectable for the OD of \acrshort{ecol26} in this preparation set, although an OD range of 1.24 to 1.79 appears to work.

To improve the \gls{conjug} protocol further, it would likely be useful to first test whether a specific growth phase of \acrlong{mbrevi} has an effect by measuring its \acrshort{OD} beforehand, because a clear trend towards fresh grown cultures is already visible for this \gls{conjug} experiment. 
To determine whether the \acrshort{OD} of \acrshort{ecol26} influences success or whether the ratio of \acrshort{ecol} cells to \gls{archa} cells does, as already described  as a critical factor in \textcite{koppCriticalVariationsConjugational2004} for the \gls{conjug} of \textit {Sorangium cellulosum}, further test series must be conducted. Here, \acrshort{OD} measurement of the used \acrlong{mbrevi} cultures is needed as well.


%%%%%%%%%%%%%%%%%%%%%%%%%%%%%%%%%%%%%%%%%%%%%%%%%%%%%%%%%%%%%%%%%%

%%%%%%%%%%%%%%%%%%%%%%%%%%%%%%%%%%%%%%%%%%%%%%%%%%%%%%%%%%%%%%%%%%

%%%%%%%%%%%%%%%%%%%%%%%%%%%%%%%%%%%%%%%%%%%%%%%%%%%%%%%%%%%%%%%%%%



%%%%%%%%%%%%%%%%%%%%%%%%%%%%%%%%%%%%%%%%%%%%%%%%%%%%%%%%%%%%%%%%%%

%%%%%%%%%%%%%%%%%%%%%%%%%%%%%%%%%%%%%%%%%%%%%%%%%%%%%%%%%%%%%%%%%%

%%%%%%%%%%%%%%%%%%%%%%%%%%%%%%%%%%%%%%%%%%%%%%%%%%%%%%%%%%%%%%%%%%


\subsection{Geography}

\subsubsection{Evaluation of immunization of cattle" as a mitigation strategy}
\label{sec:discuss_eval_immun}

As outlined in \cref{sec:reference_data} two reference data are used for the assessment of the potential impact of the mitigation strategy "immunization of cattle". On the global scale the growth rate of methane (20.9~Tg/year) for the decade 2010 to 2019 \parencite{saunoisGlobalMethaneBudget2025}. On the national scale, the emission targets of the \acrfull{gmp} to reduce the national emission about 30~\% compared to 2020 methane emissions and 40~\% to 2010 \parencite{appelhansUnterschatztesTreibhausgasMethan2025}. \acrshort{deu}, \acrshort{nzl} and \acrshort{usa} signed the \acrshort{gmp}.\\

\textbf{National scale}

\noindent For the evaluation of whether mitigation of methane production in the enteric fermentation processes of cattle can be an important missing piece in achieving the signed emission targets, \ch{CH4} emissions had to be predicted for the year 2030. As described in \cref{sec:prediction_2030}, total emissions were predicted without any additional mitigation, as well as for different success scenarios of shifting the microbiome of cattle towards less methanogenic \gls{archa}. For this purpose, predictions were performed via three linear models using different time samples as reference: the last 5, 10, and 15 years.
\Cref{fig:pred_red_sze_deu} (\acrshort{deu}), \ref{fig:pred_red_sze_nzl} (\acrshort{nzl}), and \ref{fig:pred_red_sze_usa} (\acrshort{usa}) show that the results of the 10- and 15-year models are quite similar, while the 5-year model usually differs somewhat more. This is likely due to the smaller data sample of the 5-year model, making it more susceptible to changes occurring in the most recent five years.
By examining the trends in national methane emissions for individual sectors (see \cref{fig:sectors_emission}), the linear models for each country can be evaluated.


For \acrshort{deu}, it can even be concluded that the 5-year model, predicting a lower methane emission reduction than the other two models (23.5~Tg to 26.5~Tg~\ch{CH4}), is likely the most plausible (\cref{fig:pred_red_sze_deu}). This can be explained by the fact that, in \acrshort{deu}, the energy and waste sectors have remained largely unchanged over the past six years of the dataset (2018--2024), and since emissions are already at such a low level, there is little room for further reduction (\cref{fig:sectors_emission}). From 1990 to 2018, \ch{CH4} emissions from these two sectors decreased continuously. This trend is still factored into the 10- and 15-year models, leading to predictions with a stronger decrease.
For \acrshort{nzl} (\cref{fig:pred_red_sze_nzl}) and \acrshort{usa} (\cref{fig:pred_red_sze_usa}), no such clear attribution can be made. However, the 5-year model can be described as optimistic in both cases, as it assumes that the sharper reduction in methane levels observed over the past few years will continue at this rate. The 10- and 15-year models, on the other hand, are somewhat more pessimistic, as they reflect periods of several years with stronger reductions in the past, which were subsequently followed by years of stagnant methane emissions or even slight increases.


For the three reduction scenarios tested in \cref{fig:pred_red_sze_deu}, \ref{fig:pred_red_sze_nzl}, and \ref{fig:pred_red_sze_usa} - a 10~\%, 30~\%, and 50~\% microbiome shift in the rumen towards less methanogenic \gls{archa} - it can be said that the 10~\% scenario has too little impact in all three countries to help reach the targets, if these have not already been achieved. The latter case only applies to the 10- and 15-year models of \acrshort{deu}. Therefore, further mitigation of methane emissions is needed in \acrshort{deu}, \acrshort{nzl}, and \acrshort{usa}, either to enable these countries to meet their emission targets at all, or, in the case of \acrshort{deu}, to ensure that they are met.
The 30~\% reduction scenario ensures that \acrshort{nzl} and \acrshort{deu} would meet at least one of the two reduction targets across all models, although their \acrshort{ci} might not fully achieve the emission goals. For \acrshort{nzl}, this refers to a 30~\% reduction in national methane emissions in 2030 relative to 2020, and for \acrshort{deu}, a 40~\% reduction relative to 2010. The fact that, for \acrshort{deu}, the 40~\% reduction relative to 2010 target corresponds to a higher emission level than the 30~\% reduction relative to 2020 target, whereas the opposite holds true for the other two countries, can be explained by a stronger decrease in methane emissions between 2010 and 2020 for \acrshort{deu}, while total emission values for the other two countries remained nearly stable during this period. The 50~\% reduction scenario would even be sufficient for these two countries, with all models and their \acrshort{ci} falling below both emission targets (\cref{fig:pred_red_sze_deu} and \ref{fig:pred_red_sze_nzl}).\clearpage

The situation differs somewhat for the \acrshort{usa}. In none of the scenarios do the pessimistic models (10- and 15-year) meet either of the two emission targets of the \acrshort{gmp} (see \cref{fig:pred_red_sze_usa}). The optimistic model (5-year) would come close to the 70~\% target under the 30~\% reduction scenario and fall short of it under the 50~\% reduction scenario, but would still meet the 60~\% target relative to 2010.
This difference between countries can be attributed to the differing sectoral composition of national emissions (\cref{fig:sectors_emission}). While agriculture was the main \ch{CH4} source for \acrshort{deu} and \acrshort{nzl} in 2024, in the \acrshort{usa}, energy and agriculture contribute equally, and waste, though smaller, still has a substantial impact.

Two conclusions can thus be drawn. First, if the agriculture sector, and therefore \acrshort{enterfe}, is not the dominant methane emission source for a country, addressing it alone will not be sufficient to achieve the \acrshort{gmp} targets. For example, in the \acrshort{usa}, cattle emission reductions could only become the missing piece to meet the \acrshort{gmp} emission goals if the reduction observed in the energy and waste sectors over the past five years continues at a similar rate. However, if efforts to reduce methane emissions from the energy and waste sectors were intensified, the need for cattle mitigation measures might initially be limited.
Second, if agriculture is the dominant sector, there appears to be a need for new mitigation options to meaningfully reduce methane production, as agricultural emissions did not change substantially in any of the three investigated countries (\cref{fig:sectors_emission}). \ch{CH4} emissions derived from \acrlong{enterfe} in cattle appear to be a valid target for this, provided the reduction achieved is around 30~\% or more.\\

\textbf{Global scale}

\noindent On the global scale several things gets quit obvious, if you compare the possible yearly averaged saved methane emission for the different reduction scenarios for the decade 2010-2019 (\cref{fig:saved_CH4}) with the yearly growth rate (20.9~Tg~\ch{CH4}) calculated in \citet{saunoisGlobalMethaneBudget2025} for the same decade.

First, no reduction scenario (10, 20, 30, 40, or 50~\% shift in the rumen microbiome) would have been sufficient to offset the annual growth rate. This is also true when the savings from all three countries examined are combined. Thus, neither changes in a few countries nor emission reductions in individual sectors are sufficient to stop methane accumulation in the atmosphere or even to enable a decrease in global methane levels. Nevertheless, a reduction scenario of 30 in the three countries (approx. 2.75~Tg~\ch{CH4}) alone could have offset over 10~\% of the global growth rate. Under reduction scenario 50, this figure would even be just over 20~\% at 4.5~Tg~\ch{CH4}.
Considering that climate change and rising temperatures are causing natural \ch{CH4} sources to become increasingly active, a prominent example being the thawing of permafrost \parencite{schadelPermafrostWildfireCarbon2026, saunoisGlobalMethaneBudget2025, turetskyCarbonReleaseAbrupt2020}. This means the budget for anthropogenic methane gets relatively smaller if \ch{CH4} should not accumulate in the atmosphere. Therefore, efforts to reduce methane emissions from cattle should aim for a reduction of at least around 30~\% in order to make a notable contribution to mitigate climate change. 

The second part of the \cref{fig:saved_CH4} underlines again with the calculation of \acrshort{co2_equi}, that methane reduction really can be a short-term solution to soften climate change impacts and allow so to adapt/ mitigate further. So would the saved methane emission for the reduction scenario 30 (2.75~Tg) on a time horizon of 20~years correspond to 220~Tg~\ch{CO2}, for a 100~year time horizon 75~Tg~\ch{CO2}. The \acrshort{co2_equi} for the \acrshort{gwp}20 would be roughly the same as the total \ch{CO2} emission of Spain (220.08~Tg~\ch{CO2}) for the year 2024 \parencite{europeancommission.jointresearchcentre.GHGEmissionsAll2025}.

Finally, it is noticeable in \cref{fig:saved_CH4} that the amount of methane saved increases linearly with the scenarios. In the scenarios, two values are varied: first, \acrfull{Ym}, which has a linear relationship with methane emissions (see \cref{eq:emission_total}), and second, \acrfull{DE}, which has a nonlinear mathematical relationship with methane emissions (see \cref{sec:methods-calc_saved_methane}) . Although the theoretically assumed increase in dietary efficiency leads to further reductions in \ch{CH4} emissions, this component is almost negligible because of the non-visual impact, so it has virtually no relevance to the significance of the method as a mitigation strategy.\\

\textbf{Comparison with other mitigation strategies}

\noindent Another mitigation strategy is the use of food additives, which also affect methanogenic archaea. For example, halogenated methane compounds such as bromoform inhibit the key enzyme (methyl-coenzyme M reductase) involved in methanogenesis \parencite {boadiMitigationStrategiesReduce2004, malyuginaMitigationStrategiesMethane2025, garnettGreenhousegasAbatementAustralian2024}. In this context, the addition of red algae such as \textit{Asparagopsis taxiformis} in in-vivo experiments with beef cattle resulted in an up to 80~\% reduction in methane production \parencite[Table 4]{malyuginaMitigationStrategiesMethane2025}. However, even aside from this best-case scenario, reductions of over 30~\% were regularly achieved in-vivo. To be competitive, the immunization strategy should also aim to achieve at least a 30~\% reduction.

An advantage of immunization over feed additives is that it does not require daily administration to the animals and can be easily used in grazing animals, whereas feed additives are more difficult to implement in this context because the animals are not fed continuously or there is no controlled feed composition \parencite {malyuginaMitigationStrategiesMethane2025, baca-gonzalezAreVaccinesSolution2020}.


\subsubsection{Characterization of national emission}
\label{sec:diskus_character_methan_emis}

The composition of the methane emission from the different sectors for the countries \acrshort{deu} and \acrshort{usa} fits the described global pattern in the \cref{sec:intro} and literature \parencite{appelhansUnterschatztesTreibhausgasMethan2025, mundraEmergentMethaneMitigation2024} well (\cref{fig:sectors_emission}). So they had both three dominant methane emission sources in 1990: the energy, waste and agriculture sector.
The time series (\cref{fig:sectors_emission}) shows for \acrshort{deu} clearly, methane emissions from the energy and the waste sector can already be reduced very effectively using current methods. So that in 2024 the emissions from those two sectors can hardly be reduced any further due to the low emission levels. This reduction trend is also evident for the \acrshort{usa}, albeit to a lesser extent. Although it must be noted here as a disclaimer for the strong reduction in\acrshort{deu}, the reduction in the energy sector was achieved partly by outsourcing energy production or coal extraction to other countries, which reduces emissions nationally but not globally \parencite{appelhansUnterschatztesTreibhausgasMethan2025}. At the same time, as shown in \cref{fig:sectors_emission}, the challenge of reducing emissions from the agriculture sector remains particularly striking \parencite{edelenboschReducingSectoralHardtoabate2024, intergovernmentalpanelonclimatechangeipccAgricultureForestryOther2022}.
\acrshort{nzl} also shows only minor changes in the agricultural sector. However, over the entire period under consideration from 1990 to 2024, there is no other relevant \ch{CH4} emitter besides agriculture. This clearly demonstrates that there can be notably national differences, which must be taken into account when assessing the significance and application of methods. Thus, for the United States, it is likely initially easier and more relevant to reduce emissions from energy and waste.

These national differences are also evident in the share of \acrshort{enterfe} \ch{CH4} emissions from cattle relative to total national emissions in \cref{fig:share_of_EnterFe}. Thus, \ch{CH4} from \acrlong{enterfe} processes in cattle is a major and relevant single source in all three countries. However, in the \acrshort{usa} in 2022, it accounted for only 22~\% of total emissions, whereas in \acrshort{deu} it accounted for just under 50~\% and in \acrshort{nzl} as much as 60~\%. It can, nevertheless, be noted that \acrshort{enterfe} from cattle is a good single methane emission target for the three countries examined.

When examining cattle emissions from enteric fermentation (\cref{fig:enterFe_emission_CH4}) for the three countries, it is striking that the globally described increase in cattle emissions due to rising demand for animal products \parencite{faoLivestockBalance2009, appelhansUnterschatztesTreibhausgasMethan2025} is only slightly observable. For \acrshort{deu} there is even a decrease in emissions detectable. This is due to the fact that only \acrshort{oecd} countries were considered. The described increase in demand and emissions is primarily attributable to non-\gls{annex_one} countries.

Furthermore, a comparison with the trends in cattle population numbers in \cref{fig:indic_deu}, \ref{fig:indic_nzl}, and \ref{fig:indic_usa} shows that the amount of methane emissions from cattle correlates very strongly with changes in population size. In addition, \acrshort{deu} shows that while increased production through higher-performing animals has led to an increase in \acrshort{GE}, the simultaneous reduction in the number of animals has, nevertheless, resulted in a reduction in methane emissions \parencite{Ger_Report_GHG, appelhansUnterschatztesTreibhausgasMethan2025}. However, this method of reducing population size and methane emissions cannot be continued indefinitely, and there is also the question of to what extent animal welfare is still ensured in these high-performance animals \parencite{appelhansUnterschatztesTreibhausgasMethan2025}. It can therefore be concluded that population size is the decisive factor for methane production in cattle farming.

\subsubsection{Future steps}
For the geography site next steps to improve the calculation and rating would be to check if shifting the microbiome of the rumen can have an influence on other emission sources or is cross-linked to other emission sources. One realistic option here is the manure management. From the calculation described in \textcite{ipccChapter10Emission2006} one value already seems to be  promising, the B\textsubscript{0}. B\textsubscript{0} is needed for the calculation of the \acrshort{ef} for manure management. The value represents thereby maximum methane production from the manure produced by an animal category. This could be influenced by a different ratio of methanogen to non-methanogen anaerobic microbes. But this values has to be measured, so this hypothesis, therefore, can only be tested once the first successful applications of immunization have been made.

In addition, right now, just the \gls{gross} effect was investigated off methane and its reduction in cattle. So could a counter effect be that the cow and shifted rumen microbiome produce less \ch{CH4} but more \ch{CO2}, so the shown saved carbon dioxide  in \cref{fig:saved_CH4} could be less. If the immunization strategy is in the \gls{invivo} state, this definitely needs to be checked.
Furthermore, it would be interesting to extend the calculations to other countries in order to better assess the impact of such a strategy on a global scale. Of particular interest here would be countries that are not \gls{annex_one} states, in contrast to the three examined so far, although differences in emission patterns were already evident in those countries as well (see \cref{sec:diskus_character_methan_emis})